\chapter*{Resumen}
\addcontentsline{toc}{chapter}{Resumen}

Muchos procesos de interés aplicado se registran de forma intradiaria y admiten leerse como una sucesión de curvas, una por período, cuya forma completa es la magnitud relevante para la decisión. Los modelos de series de tiempo funcionales de uso habitual, como el autorregresivo funcional (FAR), postulan una dinámica lineal con una ley de innovación fija, de modo que la distribución condicional de la curva siguiente solo se desplaza con el pasado y conserva su dispersión y su número de modos en todo estado del proceso.\\

Esta tesis propone un modelo bayesiano no paramétrico, el PSBPM-FD, que estima la distribución condicional completa de la curva siguiente dado el pasado. Las curvas se representan mediante sus componentes principales funcionales y la dinámica de cada score se modela con una mezcla probit stick-breaking de regresiones normales, extensión del Probit Stick-Breaking Process de \citet{03_02_chung2009nonparametric}, cuyos pesos dependen de los scores rezagados. Así, la forma de la ley condicional cambia con el estado reciente del proceso. La inferencia se realiza con un muestreador de Gibbs y la predictiva de los scores se transporta a la curva mediante la expansión funcional. Sobre ella se definen tres predictores puntuales (esperanza, mediana puntual y átomo más probable), tres construcciones de intervalo (equi-colas, banda simultánea y región de máxima densidad) y los comportamientos predictivos, que identifican en cada origen los desenlaces posibles de la curva siguiente y la probabilidad de cada uno.\\

El modelo se evalúa en tres escenarios simulados con heterocedasticidad condicional, cambio de régimen por umbral y multimodalidad, frente al FAR y a un Random Forest, y en el volumen horario de tráfico de la autopista I-94. En el escenario con cambio de régimen el PSBPM-FD reduce el error puntual del FAR en torno a un $21\%$ y el puntaje de Winkler en un $20\%$ a $22\%$; en los otros dos la ganancia puntual es cercana al $2\%$ y la de los intervalos al $3\%$, sostenida en la mayoría de las ventanas de prueba. La ventaja proviene de que el ancho de la banda sigue al estado del proceso, mientras que la del FAR es fija. En el tráfico, donde el tipo de día actúa como estado latente, la predictiva reparte su masa entre el perfil laborable y el de fin de semana, los comportamientos recuperan ese estado y las bandas reducen el Winkler del FAR en un $11\%$; al agregar el rezago semanal la predictiva se vuelve unimodal en los días en que esa información decide el perfil. Los resultados delimitan el alcance del modelo a series cuya ley condicional tiene más de un comportamiento: en otras cinco series reales con predictiva unimodal el PSBPM-FD queda cerca del FAR.\\

\textbf{Palabras clave:} series de tiempo funcionales, procesos de Dirichlet dependientes, probit stick-breaking, análisis de componentes principales funcional, predicción probabilística.
